\chapter{Contribution des expériences à la maîtrise des compétences}

\section{Tableau croisé Expériences $\leftrightarrow$ Blocs RNCP38152}

Les trois expériences pédagogiques analysées en Partie 2 illustrent de manière concrète ma maîtrise des compétences du référentiel RNCP38152. Le tableau suivant présente leur contribution spécifique aux 9 blocs de compétences :

\begin{table}[htbp]
\centering
\begin{tabular}{|l|c|c|c|}
\hline
\textbf{Bloc RNCP} & \textbf{FORGÉALES} & \textbf{JEUX} & \textbf{ÉVAL} \\
\hline
\textbf{BC01} & ++ & +++ & ++ \\
\hline
\textbf{BC02} & +++ & ++ & ++ \\
\hline
\textbf{BC03} & ++ & +++ & ++ \\
\hline
\textbf{BC04} & - & +++ & ++ \\
\hline
\textbf{BC05} & +++ & +++ & +++ \\
\hline
\textbf{BC06} & +++ & +++ & +++ \\
\hline
\textbf{BC07} & - & ++ & ++ \\
\hline
\textbf{BC08} & ++ & +++ & ++ \\
\hline
\textbf{BC09} & ++ & +++ & ++ \\
\hline
\end{tabular}
\caption{Tableau croisé : Contribution des expériences aux blocs RNCP38152}
\end{table}

\subsection{Légende des niveaux de contribution}

\begin{itemize}[leftmargin=*]
    \item \textbf{+++} : \textbf{Contribution majeure} - L'expérience démontre une maîtrise approfondie du bloc.
    \item \textbf{++} : \textbf{Contribution significative} - L'expérience apporte des preuves tangibles de compétence.
    \item \textbf{-} : Contribution indirecte ou non applicable.
\end{itemize}

\subsection{Analyse des contributions}

\subsubsection{Expérience FORGÉALES - Formation de formateurs à l'IA et \LaTeX{}}

\textbf{Points forts} :
\begin{itemize}[leftmargin=*]
    \item \textbf{BC01/BC03/BC07} : Conception et diffusion d'une formation numérique de formateurs avec production de ressources libres sur la Forge.
    \item \textbf{BC05/BC06} : Ingénierie de formation d'adultes (andragogie) et animation d'ateliers synchrones/asynchrones.
    \item Maîtrise des technologies émergentes (IA générative) pour l'écriture scientifique structurée.
    \item Accompagnement individualisé et collectif des enseignants et formateurs.
\end{itemize}

\subsubsection{Expérience JEUX - Pilotage de projet complexe}

\textbf{Couverture étendue des blocs} :
\begin{itemize}[leftmargin=*]
    \item \textbf{BC04/BC06} : Pilotage d'un projet complexe et collaboratif (5 disciplines, 40+ apprenants).
    \item \textbf{BC01/BC07} : Développement d'un moteur de jeu complet (outil numérique sur mesure) et diffusion sur la Forge.
    \item \textbf{BC05} : Conception et animation d'un dispositif de formation par projet.
    \item \textbf{BC03} : Communication et valorisation du projet auprès de la communauté.
\end{itemize}

\subsubsection{Expérience ÉVALUATION - Innovation en évaluation}

\textbf{Spécialisation en évaluation et métacognition} :
\begin{itemize}[leftmargin=*]
    \item \textbf{BC05/BC07} : Création d'un outil d'évaluation innovant ("Notation à l'Enveloppe").
    \item Développement de l'auto-évaluation et de la pratique réflexive chez les apprenants (BC08).
    \item Différenciation pédagogique et adaptation aux profils d'apprentissage.
    \item Modélisation d'une grille d'évaluation transférable (Ingénierie).
\end{itemize}

\section{Transfert des expériences vers l'enseignement supérieur}

\subsection{Adaptation des innovations à l'IUT}

Mon arrivée en IUT en 2025 m'a permis de commencer à transférer les approches pédagogiques développées en lycée.

\subsubsection{Adaptation de la méthode SQL}

\textbf{Contexte} : Étudiants de BUT Techniques de Commercialisation avec des niveaux hétérogènes.
\begin{itemize}[leftmargin=*]
    \item Application de la rétro-ingénierie aux systèmes d'information commerciaux.
    \item Création du besoin par l'analyse de cas concrets d'entreprises partenaires.
    \item Progression du terrain professionnel vers les concepts techniques.
\end{itemize}

\subsubsection{Projet interdisciplinaire à l'IUT}

\textbf{Projet en cours} : Projet transversal entre les départements BUT TC et Informatique.
\begin{itemize}[leftmargin=*]
    \item Développement d'applications mobiles de gestion commerciale.
    \item Collaboration entre étudiants TC (expression du besoin métier) et informatique (réalisation de la solution technique).
    \item Adaptation des méthodes agiles au rythme universitaire.
\end{itemize}

\subsubsection{Évaluation par les pairs dans le supérieur}

\textbf{Expérimentation prévue en 2025-2026} :
\begin{itemize}[leftmargin=*]
    \item Mise en place d'une évaluation croisée des projets entre groupes d'étudiants.
    \item Développement de l'esprit critique et de l'auto-évaluation.
    \item Préparation aux soutenances et à la communication professionnelle.
\end{itemize}

\subsection{Contribution à l'innovation pédagogique}

\subsubsection{Hub de ressources numériques}

\textbf{Développement en cours} :
\begin{itemize}[leftmargin=*]
    \item Adaptation des applications de la Forge aux spécificités de l'IUT.
    \item Formation de collègues aux outils numériques.
    \item Création d'une communauté de pratiques au sein du département.
\end{itemize}

\subsubsection{Recherche-action sur l'IA}

\textbf{Extension du projet de thèse} :
\begin{itemize}[leftmargin=*]
    \item Expérimentation de l'IA générative pour la personnalisation des apprentissages.
    \item Analyse de l'impact sur les étudiants à besoins particuliers en contexte universitaire.
    \item Documentation en vue d'une diffusion dans d'autres IUT.
\end{itemize}

\subsection{Rayonnement et partage}

\subsubsection{Formation des pairs}

\textbf{Actions menées et prévues} :
\begin{itemize}[leftmargin=*]
    \item Sessions de formation aux outils de la Forge des Communs.
    \item Ateliers sur la pédagogie de projet adaptée au supérieur.
    \item Partage méthodologique inter-départements.
\end{itemize}

\subsubsection{Publications et communications}

\textbf{Valorisation scientifique envisagée} :
\begin{itemize}[leftmargin=*]
    \item Rédaction d'un article sur la transposition des innovations pédagogiques du lycée à l'IUT.
    \item Proposition de communication au colloque QPES (Questions de Pédagogie dans l'Enseignement Supérieur).
    \item Contribution aux journées nationales des IUT informatique.
\end{itemize}
